\section{Lessons Learned}
\label{sec:lessons}

This section distils the transferable lessons of building and
operating the system live during the tournament, from the predictive
core, through the decision layer that turns predictions into picks, to
the engineering process around both. Each lesson was learned at the
cost of real ranking points.

\subsection{Realised form must anchor predictions for elite players}
\label{sec:lessons-form-anchor}

Every backend that survived validation was form-blind: it predicted a
player's points from price and relative team strength alone, with no
term for how the player had actually been scoring. The result was our
most consistent failure mode: under-prediction of premium
international stars such as Lionel Messi and Kylian Mbapp\'e, whom the
gradient-boosted model had never seen in its club training data and
therefore compressed toward the price-tier mean while their realised
averages ran far above it. Fed those numbers, the optimiser repeatedly
wanted to bench in-form, ceiling-priced forwards, and we overrode it
by hand. When the operators must routinely ignore the model on the
highest-variance decision in the game, the model is not doing its job.

This was not a data-availability problem: the tournament player table
carried \texttt{round\_points}, realised points per completed round,
from Matchday~1. A trailing form average was computable from day one;
it was simply never computed. The eventual correction needed two parts
working together: a leak-free lagged form feature, and retraining on
realised tournament rounds so the output scale fits the tournament.
The form feature alone, on club-only training, barely moved the
predictions; neither half is sufficient without the other.

\subsection{A designed but undeployed fix scores zero}
\label{sec:lessons-deploy}

Worse, we had anticipated the problem and built the fix, then never
turned it on. The training entry point had an \texttt{--include-wc}
flag, the code that mines realised tournament rounds into training
rows existed, and the README stated the retraining plan in plain
language. That retrain was never run in production: the deployed model
artefacts were last written on June~21 from club data only, and the
snapshot daemon re-scored the player pool every twelve hours for the
entire group stage and round of 32 against that frozen model while
four rounds of realised labels accumulated unused.

Building the retrain path and building the schedule that runs it are
two different tasks, and only the second one scores points. Likewise,
a long-lived process is a deployment target: it must log its code
version at startup and be restarted on deploy, or it will silently run
last month's code against this month's tournament. In one line: build
the loop, validate the loop, and then, separately and deliberately,
run the loop.

\subsection{Rotation risk is a first-class feature, not a fudge
factor}
\label{sec:lessons-rotation}

We modelled rotation risk as hand-set per-country multipliers in a
script: the right intuition at the wrong abstraction level. Rotation
is a per-player decision driven by manager preference, recent minutes
load, and tournament incentives such as an already-clinched group. A
proper treatment would track per-player minutes, score
manager-specific rotation propensity, detect the clinched-group signal
from standings, and output a per-player probability of starting. We
had none of this and paid for it in Matchday~3, where realised
rotation absorbed the expected gain of multiple paid transfers whose
expected-value justification was correct on its inputs. Because a
benched starter scores near zero regardless of form or matchup,
rotation is the largest unmodelled term in the system. The rule we
adopted as a margin of safety: never take an unforced transfer hit
unless the predicted gain over the next two matchdays exceeds twice
the per-hit cost.

\subsection{Captaincy is where human judgment beat the model}
\label{sec:lessons-captaincy}

Across the three group-stage matchdays our manual captain calls beat
the model's recommendation in two of the three rounds; the one loss
was a pick on which the model and the operators agreed. Both wins came
from signals the model did not consume: team news, expected manager
behaviour, and current form. No backend picked the optimal captain in
any group round. The optimiser maximises within-squad mean expected
points, but captaincy doubles a single player and rewards ceiling
rather than mean, as well as ownership differential and
fixture-specific motivation. Captain recommendations should be
advisory rather than prescriptive, and a captain objective should
score the upper quantiles of the predictive distribution, not its
mean.

\subsection{Differential strategy is portfolio construction}
\label{sec:lessons-portfolio}

We initially framed the differential question as a binary choice:
follow the template or chase differentials. The correct framing is
portfolio construction: a mix of template anchors (matching the field
on its likely captain picks, so the rank does not collapse when the
template scores), mid-ownership workhorses that accumulate the bulk of
the points, and low-ownership differential bets that create climb
opportunities when the template blanks, with the mix governed by the
current league standing. Our optimiser maximised raw expected points;
we approximated the portfolio by hand each round. A mean-variance
formulation under the budget and country constraints is the natural
formal treatment (Section~\ref{sec:future_work}).

The tournament also taught us where differentials live. The best
differential returns came from defenders and bench midfielders, while
our one attempt at a premium-forward differential returned nothing
when the player did not start. At the premium forward tier, ownership
is highly correlated with realised production; at the defender and
bench tiers the field under-attends and genuine value persists.
Template at forward, differentiate at defence and on the bench.

\subsection{Goalkeeper save value scales with opponent quality}
\label{sec:lessons-gk}

The Poisson backend priced goalkeeper save points as a flat constant,
assigning the same save expectation to a keeper behind a dominant
defence and one facing sustained pressure every match. Empirically,
expected save points scale with opponent expected goals, strongly
enough to change the goalkeeper ranking. We confirmed this as a model
bug, replaced the constant with a term scaled by opponent expected
goals, and validated the change on held-out data before shipping it
(Section~\ref{sec:negative-results} documents an earlier
mis-calibrated attempt). The general lesson: any component priced as a
constant should be audited for the covariate it silently ignores.

\subsection{Auto-substitution makes the bench an ordered queue}
\label{sec:lessons-bench}

The game's auto-substitution rule replaces a starter who plays zero
minutes with the highest-priority bench player who fits the formation,
so the bench is an ordered queue, not a set of interchangeable
backups. We never modelled the bench-order objective: the optimiser
produced a starting eleven and a bench but no ordering, which we set
by hand every round. Bench order should be a first-class solver
output, ordered by predicted points with the goalkeeper last.

\subsection{Decision-support artefacts should not need a server}
\label{sec:lessons-artefacts}

We initially planned an API service for decision review and pivoted to
a static HTML report. The static report became the highest-utility
artefact of the project: every transfer decision was made by opening
it, reading the captain board, and cross-checking against our domain
priors. The lesson generalises: produce decision artefacts viewable
without infrastructure, with Markdown for permanent records, HTML for
interactive review, and JSON for downstream consumption.

\subsection{The decision layer, not the model, was the weak layer}
\label{sec:lessons-decision-layer}

A structured mid-tournament review found no defect in the model
mathematics that had survived validation, but five distinct bugs in
the layer that turns predictions into picks: an availability discount
whose missing-history policy exempted exactly the never-played players
it existed to catch; the squad and lineup solvers silently optimising
different objective columns; a scouting bonus counted twice on one
backend's path; missing values poisoning the captain sort into
arbitrary order; and a form multiplier whose denominator grew every
round until it pinned all players to its clip floor. These are
plumbing errors, not modelling errors, and unit tests missed them
because every component met its own spec.

What was missing were end-to-end invariants on the artefact users
consume: a never-played player must not start, a captain score must be
finite, a bonus must be counted once, the budget must be within the
cap. Such invariants are cheap and would have caught most of these
bugs at introduction. Relatedly, the meaning of a missing value is a
design decision per column, not a default: three of the five bugs
trace to NaN having no declared meaning.

\subsection{Verify the data contract against a second source}
\label{sec:lessons-data-contract}

The most consequential single finding was in the data. The fantasy API
truncates each player's \texttt{round\_points} list at the last round
in which the player personally recorded anything: interior
non-appearances are zero-filled, trailing ones simply vanish, and
nothing documents this. Every consumer assumed the list length equals
the squad's completed rounds, which was false for
\TruncationShareQf{} of the player pool at the quarter-final
snapshot. The inflation lands exactly on recently-dropped players,
the population the availability discount exists to catch, and it
silently defeated the discount for all of them. The fix reconstructs
the truth from a source the API cannot truncate, the fixture table,
padding \texttt{round\_points} with zeros up to the squad's
completed-fixture count in one function shared by training and
inference. A single assertion comparing list length to completed
fixtures would have caught this in the first week; the cross-check
existed in the data the whole time, and nobody wrote it.

\subsection{Validate one component at a time, and build the harness
first}
\label{sec:lessons-validation-gate}

Twice we shipped theoretically motivated model changes as a bundle and
lost: each addition was individually defensible, and collectively they
degraded performance on every backtested round
(Section~\ref{sec:negative-results}). Simpler models outperform richer
ones on samples this small, because every added component is another
place to be slightly wrong. The decision rule we adopted as binding:
one component at a time; held-out error must improve at the position
of interest before shipping; the live tournament backtest must not
regress for at least one round; and failed attempts are documented,
because honest negatives are results too. The corollary is to build
the walk-forward validation harness first, not in week four: every
modelling decision made before ours existed was evaluated on intuition
or in-sample evidence, and most were later reversed. The harness is
the cheapest component in the repository and the only one that has
never been wrong.

\subsection{Generate every number that appears in the documentation}
\label{sec:lessons-generated-numbers}

A pre-publication audit comparing every claim in the project
documentation against the code found that hand-copied figures had
drifted in 38 places over five weeks: stale feature counts, backend
lists two versions out of date, two sections disagreeing about the
same round's realised score. Every number that was wrong had been
typed by hand; every number that was right was generated by a script
and pasted with its source cited. The lesson is mechanical rather than
moral: documentation numbers must come from script output, with the
generating script cited next to the number, and a drift audit belongs
before publication. This paper practises the rule its drafts violated:
its numeric claims are injected from generated macros rather than
typed.
